\documentclass[xcolor=table,dvipsnames,table]{beamer}
\mode<presentation>
\usetheme{boxes}
\setbeamertemplate{navigation symbols}{}
% http://www.latex-community.org/forum/viewtopic.php?f=4&t=6694
\setbeamertemplate{navigation symbols}{\raisebox{5pt}{\makebox[\paperwidth]{\hfill\makebox[10pt]{\scriptsize\insertframenumber\vspace{1ex}}}}}
%\setbeamertemplate{footline}[frame number]
\setbeamertemplate{blocks}[shadow=false]
%\setbeamercolor*{block title}{fg=structure,bg=RoyalBlue!10}
\setbeamercolor*{block title example}{fg=structure,bg=RoyalBlue!10}
%\setbeamercolor*{block title example}{fg=BrickRed,bg=Goldenrod!10}
\setbeamercolor*{block title alerted}{fg=white,bg=black}
\addtobeamertemplate{block begin}{\pgfsetfillopacity{0.8}}{\pgfsetfillopacity{1}}
%\rowcolors{0}{RoyalBlue!20}{RoyalBlue!5}
\setbeamertemplate{caption}{\raggedright\insertcaption\par}

%\DeclareGraphicsRule{*}{mps}{*}{}

\usepackage{latexsym}
\usepackage{hyperref}
\usepackage{tikz}
\usetikzlibrary{calc,shapes,arrows,shadows,shapes.callouts,shapes.arrows,chains,positioning,trees}
\usepackage{solution}
\usepackage{calc}
\usepackage{pifont}
\usepackage{algorithmic}
\usepackage{pdfcomment}
\usepackage{color}

\newcommand{\cmark}{\ding{51}}
\newcommand{\xmark}{\ding{55}}

\newcommand{\highlight}[1]{{\color{blue}{#1}}}
\newcommand{\mycite}[1]{{\color{darkgray}{\footnotesize [#1]}}}

\DeclareMathOperator*{\argmin}{arg\,min}
\DeclareMathOperator*{\argmax}{arg\,max}
\DeclareMathOperator{\sign}{sign}
\DeclareMathOperator{\cnt}{Count}

\newcounter{mycallout}

\newcommand{\callouts}[3]{%
  \stepcounter{mycallout}
  \tikz[remember picture,baseline]{\node[anchor=base,inner sep=0,outer sep=0]%
    (\themycallout) {\colorbox{#1!20}{#3}};\pause\node[overlay,rectangle callout,%
    callout relative pointer={(0cm,0.5cm)},fill=#1!20] at ($(\themycallout.south)+(-0cm,-0.7cm)$){#2};}%
    }%

\raggedright

\newcount\lecturecount
\lecturecount=0
\AtBeginLecture{%
    \advance\lecturecount by 1
    \date{}
    \begin{frame}
    \begin{center}
    \titlepage
    \ifnum\lecturecount=1
    Part \the\lecturecount: \insertlecture
    \else
    Part \the\lecturecount: \insertlecture
    \fi
    \end{center}
    \end{frame}
}

\addtobeamertemplate{block begin}{\setlength\abovedisplayskip{0pt}}

%\newcommand{\example}[1]{{\color{BrickRed!50}{#1}}}
\newcommand{\maths}[1]{{\color{RoyalBlue!50}{#1}}}
\newcommand{\reference}[1]{{\color{RoyalBlue!30}\tiny [from #1]}}
\newcommand{\koehnref}{\reference{\href{http://www.statmt.org/book}{P.Koehn SMT book slides}}}


\begin{document}

\title{\color{MidnightBlue}Natural Language Processing}

\author{Anoop Sarkar \\ {\color{RoyalBlue!70}{\href{http://anoopsarkar.github.io/nlp-class}{anoopsarkar.github.io/nlp-class}}}}
\institute{\color{BrickRed}Simon Fraser University}
%\date{}
     
{
\addtocounter{framenumber}{-1}
\begin{frame}
\begin{center}
\vspace{8mm}
\includegraphics[scale=0.35]{figures/natlang-cky-logo}
\end{center}
\titlepage
\end{frame}
}



\lecture{Probability and Language}{}

\begin{comment}
\begin{frame}
\frametitle{Announcements}
\centering
\begin{block}{Office Hours}
\smallskip
\begin{tabular}{lll}
 & Time & Location \\
\hline
Angel Chang & Fri 10:30-11:30am & TASC 8031 \\
Yasaman Etesam & Mon 2-3pm & ASB 9808 \\
            & this Fri 3:30-4:30 pm & \\
Ali Gholami & Tue 1-2 pm & ASB 9808 \\
Roya Javadi & Thu 9:30-10:30 am & ASB 9808
\end{tabular}
\end{block}

\begin{block}{Assignments}
HW0 due this Thursday
\end{block}
\end{frame}
\end{comment}

\section{Probability and Language}
\frame{\tableofcontents[currentsection]}

\begin{frame}
\frametitle{Probability and Language}
\centering
\begin{block}{Assign a probability to an input sequence}
\par Given a URL: \texttt{choosespain.com}. What is this website about?

\pause
\smallskip
\begin{tabular}{ll}
\rowcolor{MidnightBlue!50}
Input & Scoring function \\
\hline
choose spain & -8.35 \\
chooses pain & -9.88 \\
\vdots & \vdots
\end{tabular}
\end{block}

\pause
\begin{alertblock}{The Goal}
\par Find a good \textbf{scoring function} for input sequences.
\end{alertblock}
\end{frame}

\begin{frame}
\frametitle{Scoring Hypotheses in Speech Recognition}
\centering
\begin{block}{From acoustic signal to candidate transcriptions}
\begin{tabular}{ll}
\rowcolor{MidnightBlue!50}
Hypothesis & Score \\
\hline
the station signs are in deep in english & -14732 \\
the stations signs are in deep in english & -14735 \\
the station signs are in deep into english & -14739 \\
the station 's signs are in deep in english & -14740 \\
the station signs are in deep in the english & -14741 \\
the station signs are indeed in english & -14757 \\
the station 's signs are indeed in english & -14760 \\ 
the station signs are indians in english & -14790 \\
the station signs are indian in english & -14799 \\
the stations signs are indians in english & -14807 \\
the stations signs are indians and english & -14815 
\end{tabular}
\end{block}
\end{frame}

\begin{frame}[fragile]
\frametitle{Scoring Hypotheses in Machine Translation}
\centering
\begin{block}{From source language to target language candidates}
\begin{tabular}{ll}
\rowcolor{MidnightBlue!50}
Hypothesis & Score \\
\hline
we must also discuss a vision . & -29.63 \\
we must also discuss on a vision . & -31.58 \\
it is also discuss a vision . & -31.96 \\
we must discuss on greater vision . & -36.09 \\
\vdots & \vdots \\
\end{tabular}
\end{block}
\end{frame}

\begin{frame}
\frametitle{Scoring Hypotheses in Spelling Correction}
\centering
\begin{block}{Substitute spelling variants to generate hypotheses}
\begin{tabular}{p{8cm}l}
\rowcolor{MidnightBlue!50}
Hypothesis & Score \\
\hline
... stellar and versatile {\color{MidnightBlue}\textbf{acress}} whose combination of sass and glamour has defined her ... & -18920 \\
... stellar and versatile {\color{MidnightBlue}\textbf{acres}} whose combination of sass and glamour has defined her ... & -10209 \\
... stellar and versatile {\color{MidnightBlue}\textbf{actress}} whose combination of sass and glamour has defined her ... & -9801
\end{tabular}
\end{block}
\end{frame}

\begin{frame}
\frametitle{Scoring Hypotheses in Decryption}
\centering
\begin{block}{Character substitutions on ciphertext to plaintext candidates}
\begin{tabular}{ll}
\rowcolor{MidnightBlue!50}
Hypothesis & Score \\
\hline
Heopaj, zk ukq swjp pk gjks w oaynap? & -93\\
Urbcnw, mx hxd fjwc cx twxf j bnlanc? & -92\\
Wtdepy, oz jzf hlye ez vyzh l dpncpe? & -91\\
Mjtufo, ep zpv xbou up lopx b tfdsfu? & -89\\
Nkuvgp, fq aqw ycpv vq mpqy c ugetgv? & -87\\
Gdnozi, yj tjp rvio oj fijr v nzxmzo? & -86\\
Czjkve, uf pfl nrek kf befn r jvtivk? & -85\\
Yvfgra, qb lbh jnag gb xabj n frperg? & -84\\
Zwghsb, rc mci kobh hc ybck o gsqfsh? & -83\\
Byijud, te oek mqdj je adem q iushuj? & -77\\
Jgqrcl, bm wms uylr rm ilmu y qcapcr? & -76\\
Listen, do you want to know a secret? & -25\\
\end{tabular}
\end{block}
\end{frame}

\begin{frame}
\frametitle{The Goal}
\begin{itemize}[<+->]
\item Write down a \textbf{model} over sequences of words or letters.
\item \textbf{Learn} the parameters of the model from data.
\item Use the model to \textbf{predict} the probability of new sequences.
\end{itemize}
\end{frame}

\begin{frame}
\frametitle{Probabilistic Language Modeling}
\begin{itemize}[<+->]
\item Compute the probability of a sentence of sequence of words
      $P(W) = P(w_1, w_2, w_3, w_4, \ldots, w_n)$
\item Build the sequence up over time and use chain rule
  \begin{itemize}[<+->]
  \item $P(w_1, w_2, w_3, w_3, w_4, \ldots, w_n) = P(w_1) P(w_2, \ldots, w_n | w_1)$ 
  \item $P(w_1, w_2, w_3, w_3, w_4, \ldots, w_n) = P(w_1) P(w_2 | w_1) P(w_3 | w_1, w_2) \ldots P(w_n | w_1, \ldots, w_{n-1}))$ 
  \item $P(w_1, w_2, w_3, w_3, w_4, \ldots, w_n) = \prod_i P(w_i | w_1, \ldots w_{i-1})$
  \item $P(\text{it is very cold today}) = P(\text{it}) P(\text{is}|\text{it}) P(\text{very}|\text{it is}) P(\text{cold}|\text{it is very}) P(\text{today}|\text{it is very cold})$
  \end{itemize}
\item Need to compute: $P(w_i | w_1, \ldots w_{i-1})$
  \begin{itemize}[<+->]
    \item Let's count!
    \item $P(\text{cold}|\text{it is very}) = \frac{\text{Count}(\text{it is very cold})}{\text{Count}(\text{it is very})}$ 
  \end{itemize}
\end{itemize}
\end{frame}


\lecture{Quick guide to probability theory}{}

\section{Quick guide to probability theory}
\frame{\tableofcontents[currentsection]}

\begin{frame}
\frametitle{Probability: The Basics}
\begin{itemize}[<+->]
\item Sample space
\item Event space
\item Random variable
\end{itemize}
\end{frame}

\begin{frame}
\frametitle{Probability distributions}
\begin{itemize}[<+->]
\item P(X): probability of random variable X having a certain value.
\begin{itemize}[<+->]
\item P(X = killer) = 1.05e-05 
\item P(X = app) = 1.19e-05
\end{itemize}
\end{itemize}
\end{frame}


\begin{frame}
\frametitle{Joint probability}
\begin{itemize}[<+->]
\item P(X,Y): probability that X and Y each have a certain value.
\begin{itemize}[<+->]
\item Let Y stand for choice of a word
\item Let X stand for the choice of a word that occurs before Y
\item P(X = killer, Y = app) = 1.24e-10
\end{itemize}
\end{itemize}
\pause
\begin{alertblock}{Joint Probability: P(X=value AND Y=value)}
\begin{itemize}
\item Since X=value AND Y=value, the order does not matter
\item {\color{blue} P(X = killer, Y = app)} $\Leftrightarrow$ {\color{red} P(Y = app, X = killer)}
\item In both cases it is P(X,Y) = P(Y,X) = P('killer app')
\item In NLP, we often use numerical indices to express this: \\
P($W_{i-1}$ = killer, $W_{i}$ = app)
\end{itemize}
\end{alertblock}
\end{frame}

\begin{frame}
\frametitle{Joint probability}
\centering
\begin{block}{Joint probability table}
\begin{tabular}{lll}
\rowcolor{MidnightBlue!50}
$W_{i-1}$ & $W_i$ = app & P($W_{i-1}$,$W_{i}$) \\
$\langle S \rangle$ & app & 1.16e-19 \\
an & app & 1.76e-08 \\
killer & app & 1.24e-10 \\
the & app & 2.68e-07 \\
this & app & 3.74e-08 \\
your & app &  2.39e-08
\end{tabular}

\smallskip
There will be a similar table for each choice of $W_i$.
\end{block}
\begin{block}{Get P($W_i$) from P($W_{i-1}$,$W_{i}$)}
\[ P(W_i=\textsf{app}) = \sum_{x} P(W_{i-1}=x,W_i=\textsf{app}) = 1.19e-05 \]
\end{block}
\end{frame}

\begin{frame}
\frametitle{Conditional probability}
\begin{itemize}[<+->]
\item P($W_{i} \mid W_{i-1}$): probability that $W_{i}$ has a certain value after fixing value of $W_{i-1}$.
\item P($W_{i} = \textsf{app} \mid W_{i-1} = \textsf{killer}$)
\item P($W_{i} = \textsf{app} \mid W_{i-1} = \textsf{the}$)
\end{itemize}
\pause
\begin{block}{Conditional probability from Joint probability}
\[ P(W_{i} \mid W_{i-1}) = \frac{P(W_{i-1}, W_i)}{P(W_{i-1})} \]
\begin{itemize}
\item P(killer) = 1.05e-5
\item P(killer, app) = 1.24e-10
\item P(app $\mid$ killer) = 1.18e-5
\end{itemize}

\end{block}
\end{frame}

%\begin{frame}
%\frametitle{Conditional probability}
%\begin{itemize}[<+->]
%\item Conditional and joint probabilities are related:
%\[ P(X \mid Y) = \frac{P(X,Y)}{P(Y)} \]
%
%%\[ P(Y = \mbox{spain} \mid X = \mbox{choose}) &=& \frac{P(X = \mbox{choose}, Y = \mbox{spain})}{P(\mbox{spain})} \]
%%\[ = \frac{2.43 * 10^{-10}}{5.03 * 10^{-5}} \]
%%\[ = 8.78 * 10^{-5} \]
%\end{eqnarray*}
%\end{frame}
%
%\begin{frame}
%\frametitle{Bayes rule}
%\begin{itemize}[<+->]
%\item Conditional probability re-written as likelihood times prior:
%\[ P(X \mid Y) = \frac{P(Y \mid X) \times P(X)}{P(Y)} \]
%
%\begin{itemize}[<+->]
%\item $P(\mbox{choose} \mid \mbox{spain}) =
%  \frac{P(\mbox{choose} \mid \mbox{spain}) \times P(\mbox{choose})}{P(\mbox{spain})} = \frac{1.0 \times 0.001}{0.5} = 0.002$
%\end{itemize}
%\end{itemize}
%\end{frame}
%
%\begin{frame}
%\frametitle{Chain Rule}
%\begin{center}
%\begin{eqnarray}
%P(X \mid Y) &= &\frac{P(X,Y)}{P(Y)} \\
%P(Y \mid X) &= &\frac{P(Y,X)}{P(X)} \\
%P(X,Y) &= &P(Y,X) \\
%P(X \mid Y) \times P(Y) &= &P(Y \mid X) \times P(X) \\
%P(X \mid Y) &= &\frac{P(Y \mid X) \times P(X)}{P(Y)} \\
%P(X \mid Y) &= &P(Y \mid X) \times P(X)
%\end{eqnarray}
%\end{center}   
%
%\end{frame}

\begin{frame}
\frametitle{Basic Terms}
\begin{itemize}[<+->]
\item $P(e)$ -- {\em a priori} probability or just {\em prior}
\item $P(f \mid e)$ -- {\em conditional} probability. The chance of
  $f$ given $e$
\item $P(e, f)$ -- {\em joint} probability. The chance of $e$ and $f$
  both happening.
\item If $e$ and $f$ are {\em independent} then we can write $P(e,f) =
  P(e) \times P(f)$
\item If $e$ and $f$ are not {\em independent} then we can write
  $P(e,f) = P(e) \times P(f \mid e)$ \\
 $P(e,f) = P(f) \times\ ?$
\end{itemize}

\end{frame}

\begin{frame}
\frametitle{Basic Terms}
\begin{itemize}[<+->]
\item Addition of integers: \[ \sum_{i=1}^n i = 1 + 2 + 3 + \ldots +
  n \]
\item Product of integers: \[ \prod_{i=1}^n i = 1 \times 2 \times 3
  \times \ldots \times n \]
\item Factoring: \[ \sum_{i=1}^n i \times k = k + 2k + 3k + \ldots + nk
  = k \sum_{i=1}^n i \]
\item Product with constant: \[ \prod_{i=1}^n i \times k = \pause 1 k \times 2 k \ldots \times n k = \pause k^n \times \prod_{i=1}^n i \]
\end{itemize}

\end{frame}

\begin{frame}
\frametitle{Probability: Axioms}
\begin{itemize}[<+->]
\item $P$ measures total probability of a set of events
\item $P(\emptyset) = 0$
\item $P(\mbox{all events}) = 1$
\item $P(X) \leq P(Y)$ for any $X \subseteq Y$
\item $P(X) + P(Y) = P(X \cup Y)$ provided that $X \cap Y = \emptyset$
\end{itemize}
\end{frame}

\begin{frame}
\frametitle{Probability Axioms}
\begin{itemize}[<+->]
\item All events sum to 1: \[ \sum_e P(e) = 1 \]
\item Marginal probability $P(f)$: \[ P(f) = \sum_e P(e, f)  \]
\item Conditional probability: \[ \sum_e P(e \mid f) = \pause \sum_e \frac{P(e,f)}{P(f)} = \pause \frac{1}{P(f)} \sum_e P(e,f) = \pause 1 \]
\item Computing $P(f)$ from axioms: \[ P(f) = \sum_e P(e) \times
  P(f \mid e) \]
\end{itemize}

\end{frame}

\begin{frame}
\frametitle{Probability: The Chain Rule}
\begin{itemize}[<+->]
\item $P(a,b,c,d \mid e)$
\item We can simplify this using the Chain Rule:
\item $P(a,b,c,d \mid e) = $\\
$P(d \mid e) \cdot P(c \mid d,e) \cdot P(b \mid c,d,e) \cdot P(a \mid b,c,d,e)$ 
\item To see why this is possible, recall that \( P(X \mid Y) = \frac{p(X,Y)}{p(Y)} \)
\begin{itemize}
\item $\frac{ p(a,b,c,d,e) }{ p(e) } = \frac{p(d,e)}{p(e)} \cdot \frac{p(c,d,e)}{p(d,e)} \cdot \frac{p(b,c,d,e)}{p(c,d,e)} \cdot \frac{p(a,b,c,d,e)}{p(b,c,d,e)} $
\end{itemize}
\item We can approximate the probability by removing some random variables from the context. For example, we can keep at most two variables to get:
\[ P(a,b,c,d \mid e) \approx P(d \mid e) \cdot P(c \mid d, e) \cdot P(b \mid c, e) \cdot P(a \mid b, e) \]
\end{itemize}

\end{frame}

\begin{frame}
\frametitle{Probability: The Chain Rule}
\begin{itemize}[<+->]
\item $P(e_1, e_2, \ldots, e_n) = P(e_1) \times P(e_2 \mid
  e_1) \times P(e_3 \mid e_1, e_2) \ldots $ 
\[ P(e_1, e_2, \ldots, e_n) = \prod^n_{i=1} P(e_i \mid e_{i-1},
  e_{i-2}, \ldots, e_1) \]
\item In NLP, we call:
	\begin{itemize}[<+->]
	\item $P(e_i)$: unigram probability 
	\item $P(e_i \mid e_{i-1})$: bigram probability
	\item $P(e_i \mid e_{i-1}, e_{i-2})$: trigram probability
	\item $P(e_i \mid e_{i-1}, e_{i-2}, \ldots, e_{i-{(n-1)}})$: n-gram probability
	\end{itemize}
\end{itemize}

\end{frame}

\begin{frame}
\frametitle{Probability: Random Variables and Events}
\begin{itemize}[<+->]
\item What is $y$ in $P(y)$ ?
\item Shorthand for value assigned to a random variable $Y$, e.g. $Y = y$
\item $y$ is an element of some implicit {\bf event space}: ${\cal E}$
\end{itemize}

\end{frame}

\begin{frame}
\frametitle{Probability: Random Variables and Events}
\begin{itemize}[<+->]
\item The {\em marginal probability} $P(y)$ can be computed from $P(x,y)$ as
  follows:
\[ P(y) = \sum_{x \in {\cal E}} P(x,y) \]
\item Finding the value that maximizes the probability value:
\[ \hat{x} = \arg\max_{x \in {\cal E}} P(x) \]
\end{itemize}

\end{frame}

\begin{frame}
\frametitle{Log Probability Arithmetic}
\begin{itemize}[<+->]
\item Practical problem with tiny $P(e)$ numbers: underflow
\item One solution is to use log probabilities:
\begin{eqnarray}
\log(P(e)) &=& \log(p_1 \times p_2 \times \ldots \times p_n) \nonumber
\\
&=& \log(p_1) + \log(p_2) + \ldots + \log(p_n) \nonumber
\end{eqnarray}
\item Note that: \[ x = \exp(\log(x)) \]
\item Also more efficient: addition instead of multiplication
\end{itemize}

\end{frame}

\begin{frame}
\frametitle{Log Probability Arithmetic}
\begin{center}
\begin{tabular}{|c|c|}
\rowcolor{MidnightBlue!50}
$p$ & $\log(p)$ \\
\hline
$0.0$ & $- \infty$ \\
$0.1$ & $-3.32$ \\
$0.2$ & $-2.32$ \\
$0.3$ & $-1.74$ \\
$0.4$ & $-1.32$ \\
$0.5$ & $-1.00$ \\
$0.6$ & $-0.74$ \\
$0.7$ & $-0.51$ \\
$0.8$ & $-0.32$ \\
$0.9$ & $-0.15$ \\
$1.0$ & $0.00$ \\
\hline
\end{tabular}
\end{center}

\end{frame}

\begin{frame}
\frametitle{Log Probability Arithmetic}
\begin{itemize}[<+->]
\item So: $(0.5 \times 0.5 \times \ldots 0.5) = (0.5)^n$ might get too
  small but $(- 1 - 1 - 1 - 1) = - n$ is manageable
\item Another useful fact when writing code \\
($\log_2$ is {\em log to the base 2}):
\[ \log_2(x) = \frac{\log_{10}(x)}{\log_{10}(2)} \]
\end{itemize}

\end{frame}

\lecture{Entropy and Information Theory}{}

\section{Entropy and Information Theory}
\frame{\tableofcontents[currentsection]}

\begin{frame}
\frametitle{Information Theory}
\begin{itemize}[<+->]
\item Information theory is the use of probability theory to quantify
  and measure ``information''.
\item Consider the task of efficiently sending a message. Sender Alice
  wants to send several messages to Receiver Bob. Alice wants to do
  this as efficiently as possible.
\item Let's say that Alice is sending a message where the entire
  message is just one character {\em a}, e.g. {\em aaaa$\ldots$}. In
  this case we can save space by simply sending the length of the
  message and the single character.
\end{itemize}

\end{frame}

\begin{frame}
\frametitle{Information Theory}
\begin{itemize}[<+->]
\item Now let's say that Alice is sending a completely random signal to
  Bob. If it is random then we cannot exploit anything in the message
  to compress it any further.
\item The {\em expected} number of bits it takes to transmit
  some infinite set of messages is what is called entropy. 
\item This formulation of entropy by Claude Shannon was adapted from
  thermodynamics, converting information into a quantity that can be measured.
\item Information theory is built around this notion of message
  compression as a way to evaluate the amount of information. 
\end{itemize}

\end{frame}

\begin{frame}[fragile]
\frametitle{Expectation}
\begin{itemize}[<+->]
\item For a probability distribution $p$
\item \textbf{Expectation} with respect to $p$ is a weighted average:
\begin{eqnarray*}
 E_p[x] &=& \frac{x_1 \cdot p(x_1) + x_2 \cdot p(x_2) + \ldots + x_n p(x_n)}{p(x_1) + p(x_2) + \ldots + p(x_n)} \\
  &=& x_1 \cdot p(x_1) + x_2 \cdot p(x_2) + \ldots + x_n p(x_n) \\
  &=& \sum_{x \in {\cal E}} x \cdot p(x) 
\end{eqnarray*}
\item Example: for a six-sided die the expectation is:
\[ E_p[x] = 1 \cdot \frac{1}{6} + 2 \cdot \frac{1}{6} + \ldots + 6 \cdot \frac{1}{6} = 3.5 \]
\end{itemize}

\end{frame}

\begin{frame}
\frametitle{Entropy}
\begin{itemize}[<+->]
\item For a probability distribution $p$
\item \textbf{Entropy} of $p$ is: 
\[ H(p) = - \sum_{x \in {\cal E}} p(x) \cdot \log_2 p(x) \]
\item Any base can be used for the $\log$, but base $2$ means
  that entropy is measured in bits.
\item What is the {\it expected} number of bits with respect to $p$:
\[ - E_p [ \log_2 p(x) ] = H(p) \]
\item Entropy answers the question: {\it What is the expected
  number of bits needed to transmit messages from event 
  space ${\cal E}$, where $p(x)$ defines the probability of observing $x$?} 
\end{itemize}

\end{frame}

\begin{comment}
	
\begin{frame}
\frametitle{Entropy}
\begin{itemize}[<+->]
\item Alice wants to bet on a horse race. She has to send a message to
  her bookie Bob to tell him which horse to bet on.
\item There are 8 horses. One encoding scheme for the messages is to
  use a number for each horse. So in bits this would be $001, 010,
  \ldots$\\
(lower bound on message length = 3 bits in this encoding scheme) 
\item Can we do better?
\end{itemize}

\end{frame}

\begin{frame}
\frametitle{Entropy}
\begin{center}
\begin{tabular}{|cc|cc|}
\hline
Horse 1 & $\frac{1}{2}$ & Horse 5 & $\frac{1}{64}$ \\
\hline
Horse 2 & $\frac{1}{4}$ & Horse 6 & $\frac{1}{64}$ \\
\hline
Horse 3 & $\frac{1}{8}$ & Horse 7 & $\frac{1}{64}$ \\
\hline
Horse 4 & $\frac{1}{16}$ & Horse 8 & $\frac{1}{64}$ \\
\hline
\end{tabular}
\end{center}
\begin{itemize}[<+->]
\item If we know how likely we are to bet on each horse, say based on
  the horse's probability of winning, then we can do better.
\item Let $p$ be the probability distribution given in the
table above. The entropy of $p$ is $H(p)$
\end{itemize}

\end{frame}

\begin{frame}
\frametitle{Entropy}
\begin{eqnarray*}
\lefteqn{ H(p) = } \nonumber \\
&=& - \sum^{8}_{i=1} p(i) \log_2\ p(i) \nonumber \\
&=& - \left( \frac{1}{2} \log_2 \frac{1}{2} + \frac{1}{4} \log_2 \frac{1}{4} +
\frac{1}{8} \log_2 \frac{1}{8} + \frac{1}{16} \log_2 \frac{1}{16} +
4 (\frac{1}{64} \log_2 \frac{1}{64}) \right) \nonumber \\
&=& - \left( \frac{1}{2} \times -1 + \frac{1}{4} \times -2 +
\frac{1}{8} \times -3 + \frac{1}{16} \times -4 +
4 (\frac{1}{64} \times -6 ) \right) \nonumber \\
&=& - \left( {\color{red}-\frac{1}{2} - \frac{1}{2}} {\color{blue} - \frac{3}{8} - \frac{1}{4} - \frac{3}{8}} \right) \nonumber \\
&=& 2\ \textit{bits}
\end{eqnarray*}

\begin{itemize}[<+->]
\item What is the entropy when the horses are equally likely to win?
\[ H(\textit{uniform distribution}) = - 8 (\frac{1}{8} \times -3) = 3 \textit{ bits} \]
\end{itemize}

\end{frame}

\begin{frame}
\frametitle{Entropy}
\begin{itemize}[<+->]
\item e.g., most likely horse gets code $0$, next most likely gets $10$, and then $110, 1110, \ldots$ \\ {\small many possible coding schemes, this is a simple code to illustrate number of bits needed for a large number of messages \ldots}
\item Assume there are $320$ messages (one for each race): \\
code $0$ occurs 160 times, code $10$ occurs 80 times, code $110$ occurs 40 times, code $1110$ occurs 20 times, code $11110$ occurs 5 times.
\item Total number of bits for all messages: 160*len(0) + 80*len(10) + 40*len(110) + 20*len(1110) + 5*len(11110) 
\item Number of bits: 160*1 + 80*2 + 40*3 + 20*4 + 5*5 = 545
\item Total number of bits per message (per race): $\frac{545}{320} \approx$ 1.7 bits (always less than 2 bits)
\end{itemize}

\end{frame}
\end{comment}

\begin{frame}
\frametitle{Perplexity}
\begin{itemize}[<+->]
\item The value $2^{H(p)}$ is called the {\bf perplexity} of a distribution $p$
\item Perplexity is the weighted average number of choices a random
  variable has to make.
\item Choosing between 8 equally likely options (H=3) is $2^3 = 8$.
\end{itemize}

\end{frame}

\begin{frame}
\frametitle{Relative Entropy}
\begin{itemize}[<+->]
\item We often wish to determine the divergence of a distribution $q$ from another distribution $p$
\item Let's say $q$ is the estimate and $p$ is the true probability
\item We define the {\em divergence} from $q$ to $p$ as the {\bf relative entropy}: written as $D(p \| q)$
\[ D(p \| q) = - \sum_{x \in {\cal E}} p(x) \log_2 \frac{q(x)}{p(x)} \]
\item Note that
\[ D(p \| q) = - E_{ p(x) } \left[ \log_2 \frac{q(x)}{p(x)} \right] \]
\item The relative entropy is also called the {\em Kullback-Leibler divergence}.
\end{itemize}

\end{frame}

\begin{frame}
\frametitle{Cross Entropy and Relative Entropy}
\begin{itemize}[<+->]
\item The {\bf relative entropy} can be written as the sum of two terms:
\begin{eqnarray*}
D(p \| q) & = & - \sum_{x \in {\cal E}} p(x) \log_2 \frac{q(x)}{p(x)} \\
& = & - \sum_{x} p(x) \log_2 q(x) + \sum_{x} p(x) \log_2 p(x) 
\end{eqnarray*}
\item We know that $H(p) = - \sum_{x} p(x) \log_2 p(x)$
\item Similarly define $H(p,q) = - \sum_{x} p(x) \log_2 q(x)$
\[
\begin{array}{l@{}c@{}l@{}c@{}l@{}}
D(p \| q) & = & H(p,q) & - & H(p) \\
\textbf{relative entropy}(p,q) & = & \textbf{cross entropy}(p,q) & - & \textbf{entropy}(p) 
\end{array}
\]
\item The term $H(p,q)$ is called the {\bf cross entropy}.
\end{itemize}

\end{frame}

\begin{frame}
\frametitle{Cross Entropy and Relative Entropy}
\begin{itemize}[<+->]
\item $H(p,q) \geq H(p)$ always.
\item $D(p \| q) \geq 0$ always, and $D(p \| q) = 0$ iff $p = q$
\item $D(p \| q)$ is not a true distance: 
  \begin{itemize}
  \item It is asymmetric: $D(p \| q) \neq D(q \| p)$, 
  \item It does not obey the triangle inequality: $D(p \| q) \nleq D(p \| r) + D(r \| q)$
  \end{itemize}
%\item Pinsker's inequality (sup is the lowest upper bound): 
%\[ \sqrt{ \frac{ D(p \| q) }{ 2 } } \geq \sup \{ | p(x) - q(x) | \} \]
\end{itemize}

\end{frame}

\begin{frame}
\frametitle{Conditional Entropy and Mutual Information}
\begin{itemize}[<+->]
\item {\it Entropy} of a random variable $X$:
\[ H(X) = - \sum_{x \in {\cal E}} p(x) \log_2 p(x) \]
\item {\it Conditional Entropy} between two random variables $X$ and $Y$:
\[ H(X \mid Y) = - \sum_{x,y \in {\cal E}} p(x,y) \log_2 p(x
\mid y) \]
\item {\it Mutual Information} between two random variables $X$ and $Y$:
\[ I(X;Y) = D(p(x,y) \| p(x)p(y)) = \sum_x \sum_y p(x,y) \log_2
\frac{p(x,y)}{p(x)p(y)} \]
\end{itemize}

\end{frame}

\begin{comment}
\section*{Appendix: Computing with Log Probability}

\begin{frame}[fragile]
\frametitle{Log Probability Arithmetic}
\begin{center}
\begin{verbatim}
// In Python: use numpy.logaddexp2(x1,x2) for base 2
// computes log(a + b) given log(a) and log(b)
double logadd(double lna, double lnb)
{
    if (lna == 1.0) return lnb;
    if (lnb == 1.0) return lna;
    double diff = lna - lnb;
    // 500 is log of large constant
    if (diff < 500.0) 
        return log1p(exp(diff)) + lnb;
        // log1p(x) computes log(1+x)
    else
        return lna;
}
\end{verbatim} 
\end{center}
\end{frame}
\end{comment}



\section*{Acknowledgements}

\begin{frame}
\centering
\begin{alertblock}{Acknowledgements}
Many slides borrowed or inspired from lecture notes by Michael Collins, Chris Dyer, Kevin Knight, Philipp Koehn, Adam Lopez, Graham Neubig and Luke Zettlemoyer from their NLP course materials. 

\bigskip

All mistakes are my own.

\bigskip

A big thank you to all the students who read through these notes and helped me improve them.

\end{alertblock}
\end{frame}



\end{document}

